\section{Conclusion}\label{sec:conclusion}

Client-side protection cannot win against an attacker who owns the machine, and AI-assisted
tooling shortens the time it takes to lose. SA-HEACS accepts this and keeps the code that
matters off the machine: a zero-IP client shell, attested and encrypted calls to segmented
server clusters, and an orchestrator that uses locality routing, micro-batching and prefetch
to keep remote execution affordable. A promotion and demotion pipeline moves each feature to
the client once its secrecy no longer pays for its serving cost.

The design eliminates static and dynamic recovery of cloud-only code. It does not eliminate
functional cloning through the API, which becomes the main residual risk and calls for rate
limiting, query monitoring, output shaping and watermarking. Future work should evaluate the
design on a real application with production traffic, run an AI-assisted red team against
the client shell, and measure extraction cost against the deployed controls.

\section*{Code availability}
\addcontentsline{toc}{section}{Code availability}
All programs are in \code{paper/code/} and listed in \cref{app:code}. They run with
Python~3.11 or later after \code{pip install -r requirements.txt}; \code{run\_all.py} runs
every program, saves its output to \code{code/outputs/}, and regenerates the CSV files
behind \cref{fig:batching,fig:prefetch}.
